\begin{chapterbox}
    \chapter{Stürmische Rätselnacht in der Taverne (2019 bis 2020)}
    \label{Stürmische Rätselnacht in der Taverne (2019 bis 2020)}

    \az{Jahr 66}

    Während einer Sturmnacht verbleiben viele Andori im "Trunkenen Troll". Ein betrunkener Bauer erzählt Schauergeschichten von Skralreigen. Mysteriöse Steintafeln mit unbekannten Schriftzeichen geben Gilda Rätsel auf. Die Wassermagierin Jarid und der Feuerkrieger Trieest aus dem fernen Danwar treffen einen blinden Seher.
\end{chapterbox}



\section{Die Sage vom Butterbrotbären}

\az{Jahr 66}

Es war früh am Morgen (oder noch spät in der Nacht?) als Gilda dabei war, ihr berühmtes andorisches Tavernenbrot mit feiner Butter (von ihrer Ziege Molli) zu backen. Ein ganzes Brett Butterbrote stand bereits auf den Fenstersims und verbreitete seinen frischen Brotgeruch in der gesamten Backküche. Gerade drehte Gilda sich um, um ein neues Tablett zu holen, da vernahm sie ein lautes Krachen hinter sich. Als sie sich umblickte, lag das erste Brett am Boden und die Brötchen waren nirgendwo mehr zu sehen. Geschwind eilte sie zum Fenster, um sich nach dem Dieb umzusehen – doch in der Dunkelheit war ohnehin kaum etwas zu erblicken. Als Gilda sich bückte, um wenigstens das gefallene Tablett wieder aufzuheben, fand sie jedoch darunter ein versiegeltes Pergament. Ein Brief? Schnell rollte Gilda das Papier auf – Grossvater Erloth hatte ihr schliesslich nicht umsonst Lesen und Schreiben beigebracht – und versuchte zu entziffern, was darauf stand. Doch die meisten Symbole auf dem Dokument waren ihr völlig unbekannt. Einige Kritzeleien sahen aus, als hätte jemand versucht, ein Tier zu zeichnen. Mysteriös, die ganze Sache.

Eine Woche darauf war Gilda erneut dabei, Butterbrote zu backen (und jawohl, in der Andor-Welt werden Butterbrote als solche gebacken, nicht erst später geschmiert), als sie erneut hinter sich ein Krachen vernahm. Die Brote zu retten vermochte sie nicht mehr, dafür war es schon wieder zu spät. Aber diesmal konnte Gilda einen Blick auf den Übeltäter werfen, wie er aus dem Lichtschein ihres Fensters in die Dunkelheit des Südlichen Waldes verschwand – und es war kein Mensch, vielmehr ein pelziges Ungetüm, welches geschwind, wenn auch ungeschickt, alle Butterbrote in sein Maul gepackt hatte und damit davongerannt war. Wie schon beim ersten Mal fand Gilda auf dem zurückgelassenen Butterbrotbrett ein Schriftstück. Dieses hier war allerdings eine Schiefertafel, in welche jemand – oder etwas – krakelig einige Zeichen hineingekratzt hatte. Mit etwas Fantasie könnte man darin ein Strichmännlein erkennen. Was das wohl zu bedeuten hatte?

Die nächste Woche war Gilda gewappnet. Schon wie zuvor backte sie ihre berühmten Tavernenbutterbrote und stellte ein Brett voller frisch gebackener Backwaren auf ihren Fenstersims zum Auskühlen. Doch als sie dem Fenster den Rücken zuwandte, war sie auf das Krachen gefasst. Blitzschnell fuhr sie herum, sprang behände über den ganzen Haufen von Schiefertafeln, den der Langfinger diesmal zurückgelassen hatte, und aus dem Fenster hinaus. Gilda rannte hinter dem diebischen Riesen hinterher – aha, ein Bär war das! Selbiger raste schwerfällig davon, als wäre der Teufel hinter ihm her, und verschwand im Dickicht des Südlichen Waldes. Doch Gilda gab noch nicht auf. Einst hatte sie sich ein wenig Fährtenkunde von Fenn, dem Fährtenleser, beibringen lassen, und das konnte sie jetzt nutzen. Wobei die Schneise, welche das braunfellige Ungetüm hinter sich gelassen hatte, selbst für absolute Anfänger in der Kunst des Nachstellens unglaublich offensichtlich gewesen wäre.

Der Spur der Verwüstung durch Gebüsch und Gesträuche folgend, erreichte Gilda nach einem Marsch von einigen Minuten eine kleine Hütte in einer vom Mondlicht beschienenen Waldlichtung. Das war seltsam, sie hatte nicht gedacht, dass sich hier eine befände – und sonst kannte sie sich eigentlich gut im Südlichen Wald aus. Vorsichtig näherte sie sich dem Gemäuer. Der Geruch nach frisch gebackenem Butterbrot – ihrem Brot! – lag in der noch kühlen Waldluft. Das erste Vogelgezwitscher kündigte den baldigen Sonnenaufgang an. Und Gilda trat an das Fenster des kleinen Häuschens und guckte hinein.

Das Innere der Hütte war von einem warmen Kaminfeuer erleuchtet und sah unglaublich gemütlich aus. Gilda erhaschte einen Blick auf eine grosse Stube, in deren Holzwände zahlreiche Zeichen hineingeschnitzt worden waren, ebenjene Zeichen, die sich auch auf dem Pergament und den Schiefertafeln befanden, die das butterbrotestehlende Wesen zurückgelassen hatte. Am Boden erkannte sie einen waschechten tulgorischen Teppich – wertvolle Ware war das! In einem grossen Sessel vor dem Feuer sass eine Gestalt, Gilda den Rücken zugewandt, und genoss augenscheinlich eine Tasse dampfenden Gebräus.

Gilda lehnte sich ein ganz vorsichtig ein wenig weiter vor, um zu sehen, wer denn da in dieser Hütte sass. Da knackte es unter ihren Schuhen – ein trockener Ast, das war ja auch zu erwarten gewesen – und die Gestalt im Sessel schreckte auf. Rasch zog Gilda sich zurück und rannte los zur Taverne, da krachte auch schon die Tür zur geheimnisvollen Hütte auf. Im hellen Schein aus der Hütte erblickte Gilda den diebischen Bären, welcher aus der Tür getreten war und nun bedrohlich brummelnd schnell auf sie zuhetzte. Geschwind drehte sie sich um und hetzte zurück zur sicheren Taverne, durch Unterholz und Gestrüpp, in der Hoffnung, es würde ihren massigen Verfolger etwas abbremsen.

Dem war auch so: Als Gilda die sichere Pforte der Taverne erreichte, hatte sich der Bär erst gerade aus einem Dornenstrauch am Ende des Südlichen Walds befreit. Und wie wir alle wissen, dürfen Bären die Taverne zum Trunkenen Troll ja nicht betreten. Somit blieb dem Bären nichts übrig, als sich nach einem letzten Blick auf das Gebäude langsam umzudrehen und zurück in den Wald zu trotten.

Gilda beobachtete vom Fenster aus, wie das Tier das Weite suchte, während ihr Atem sich langsam wieder beruhigte. In den Tagen, die da kommen, versuchte sie ein-, zweimal im Tageslicht, die mysteriöse Hütte wiederzufinden – erfolglos, die Spuren des Bären führten ins Nichts. Und natürlich berichtete sie in der Taverne von dieser Begegnung, aber mehr ausser Orfens spöttische These, dass sie vielleicht etwas über den Durst getrunken hatte, konnte sie nicht hervorbringen. Nicht einmal bei den Bewahrern vom Baum der Lieder hatte man von dieser seltsamen Hütte gehört. Und nicht einmal Eara konnte den vom Bären zurückgelassenen Schiefertafeln Informationen entlocken, und sie hatte es weiss Mutter Natur versucht.

Seitdem stellte Gilda ihre berühmten Butterbrote jedenfalls nicht mehr frisch gebacken auf ihre Fensterbank. Bloss von Zeit zu Zeit liess sie einige wenige am Rande des Südlichen Waldes liegen. Wenn sie später wiederkam, würden die Brote bereits weg sein, das wusste sie, und manchmal fand sie sogar einige weitere unerzifferbaren Schriftstücke an deren Stelle. Aber den Butterbrotbären sah sie nie mehr – und vielleicht waren es auch einfach irgendwelche Waldgeister, die die Brote jeweils wegschnappten und sich einen Streich mit ihr erlaubten.

Hier endet die Sage vom Butterbrotbären. Gilda hofft immer noch darauf, dass eines Tages jemand die Taverne besucht, der die rätselhaften Botschaften auf den Tafeln entziffern kann.



\newpage
\section{Der Skralreigen}


Ich hab sie gesehen! Ein Skralreigen! Ich schwör’s euch, ich hab den gesehen! Weit draussen im Rietlande war ich, alleine der Narne entlangschlendernd auf dem Nachhauseweg durch das helle Mondlicht – ich muss euch ja gar nicht erzählen, wie wild die Party in der guten alten Taverne diese Nacht war, ihr wart ja selbst alle dabei, ihr elenden Halunken! Also ging ich da diesem Fluss entlang, sturzbetrunken natürlich, und da erblickte ich im Licht des roten Mondes, dass das Graue Gebirge sich bewegte!

Zunächst dacht’ ich natürlich, ich spinne! Ich mein’, das Graue Gebirge steht da schon, seitdem Mutter Natur es als Keil zwischen diese guten Lande und das düstere Krahd getrieben hatte, und einen solchen Keil wird weder wanken noch weichen, egal, was geschieht, so sagten’s diese eitlen Gesandten der Bewahrer jedenfalls, wo sie unser Dörfchen besuchten, als ich noch ein kleiner Hosenscheisser war, und uns allen Wissen und Tugend übermitteln wollten. Diese armen Kerle, ihre Saat hat bei uns Haudegen kaum gefruchtet. Aber ich schweif’ ja mal wieder grandios ab!

Ach ja, das Graue Gebirge! Eben, ich sah mit meinen eig'nen Augen, wie der Gipfel des Skralbergs sich vor mir sich etwas in die Höhe erhob und wieder niedersenkte. Das haute mich glatt von den Socken, und so setzte ich mich auf meinen gepflegten Allerwertesten und staunte die riesigen Gesteinmassen vor mir an, die immer noch herumschwankten.

Nein, es war schon keine grosse Bewegung, aber dennoch deutlich wahrnehmbar.

Halt’s Maul, Mard, willst du mich etwa einen Lügner nennen?! Du bist mehr Trunkenbold, als ich es je sein werd’!

Jedenfalls näherte ich mich dem wogenden Grauen Gebirge, und da, an einer Flussböschung, erkannte ich ein Feuerchen. Und ihr kennt mich, ich bin immer zu haben für ein Feuerchen, am Feuerchen tun Gesellschaften ihr Innerstes kund, nichts Schlechtes hat je mit ’nem Feuerchen begonnen, bla bla bla, da nähere ich mich also diesem Feuerchen und den Personen, die im Kreis drum rumsitzen.

Gilda, Schätzchen! Haste noch ’n bissle vom Rachenputzer der Schildzwerge übrig? Ein halbes Mass gerne!

Also, diese Gestalten, die da im Kreis um das Feuerchen rumsitzen, ich näher’ mich denen, und mir bleibt fast die Spucke im Mund stecken, als ich erkenn’, was die sind: Das waren Skrale! Waschechte Skrale, einer neben dem anderen schön brav um dieses grosse Feuerchen rumsitzend. Nicht nur die Krieger, die unsereins des Nachts terrorisieren kommen, nein, auch ihre Frauchen und Kindlein, und, behüte mich, selbst eine grausige Skralhexe hatte sich an der Stelle breitgemacht.

Nun, zu diesem Zeitpunkt hatt’ ich natürlich schon mit meinem bescheidenen kleinen Leben abgeschlossen. Eine ganze Skralhorde, das überlebten nicht einmal die prächtigen Helden von Andor! Nicht persönlich gemeint, Eara, ich mag dich, aber das waren bestimmt zehn, zwölf, nein, einhundert grausame Skrale! Da hättest auch du nichts mehr ausrichten können.

Die blickten mich alle an, als wär’ ich ein Geist. Und das Wunder war, die haben mich nicht angegriffen, nicht mal angefaucht. Stattdessen klopfte so ein Wichtigtuer mit seinem Stab auf den Boden – was meinst du, Eara? Wen kümmert das schon, wie der Stab ausschaute, wie so 'n Zauberstab jedenfalls, mit so’n Paar Klunkerchen dran, auf die der geizige Garz ganz schön anspringen würde. Also ja, der klopft mit seinem Stab auf den Boden, und die Skralhexe, gebeugt und bucklig überragt’ die mich immer noch mindestens drei Köpf’, also, diese Skralhexe, die gebietet mir, näher zu treten.

Kann natürlich nicht Andorisch sprechen, das ungebildete Pack, aber mit ihren Armen und einigen grossen Gesten konnt’ dieses Ungetüm von ’ner Hexe sich schon verständlich machen. Ich hab’ natürlich schon die Gerüchte gehört über die Skralreigen, die des Vollmonds am Skralberg tagen, um Dunkle Rituale zu vollbringen, und ich wusst’ weit hinten in meinem sturzbesoffenen Hinterkopp sicher auch, dass sie in jenen Nächten ums Verrecken kein Blut vergiessen wollen, weil, wenn man den Geschichten Glauben schenkt, so kriegen die Skrale in diesen Vollmondnächten ihren dreckigen Nachwuchs, und ich mein’, wenn man diese Zeit der Geburt mit Kriegereien verbringt, so wird der Nachwuchs schwach und verkümmert, weil man von seiner Kriegskunst stiehlt, bla bla – aber hey! Ich weiss’ ja offensichtlich doch noch’n bissle was davon, was diese Bewahrer uns beibringen wollten!

Ja, natürlich ist das Quark mit Sosse, diese Skrale glauben doch absoluten Blödsinn, aber jedenfalls würden sie diese Nacht niemanden angreifen, auch nicht mich, ich war also vergleichsweise sicher, solange der Vollmond da stand.

Trotzdem schlotterte ich wie Eschenlaub, als ich mich denen näherte. Und angewidert war ich auch. Aber auch betrunken bis zum Gehtnimmer, also dacht’ ich, was solls, und setz’ mich zu diesen Kerlen ans Feuer. Die sitzen da rum und murmeln ’n bissle was, so’n kehliger Singsang halt, die alte Hexe am lautesten. Einige blinzelten mir entgegen und ich blinzelte zurück und einige starrten mich nur aus dem Augenwinkel an und nein wirklich, was mach’ ich eigentlich hier?! Ich konnt’ di ganze Zeit nur daran denken, was die Skrale dem armen Gerdan und seiner Familie angetan hatten. Da hatt’ ich mich schon anders besonnen und woll’t aufstehen und mich wieder vom Acker machen, Dunkle Hexerei war das, was die da taten, Dunkle Hexerei, das sag’ ich euch, aber als ich Anstalten machte, aufzustehen, packte mich einer dieser sitzenden Grobiane mit einer seiner dreckigen Tatzen an der Schulter und zog mich wieder nach unten! Die wollten nicht, dass ich abhaute! Da fiel mir das Herz schon ’n bissle in die Hose. Ich wusst’ ja nicht, was für’n Ritual so’n Skralreigen abhält und ob die dafür Opfer brauchen oder so. War mir aber sicher, dass sie mich nicht einfach nur in den Kreis eingeladen hatten, um mir ihre Kultur oder so zu präsentieren, die geben ja nicht viel auf Menschen, ausser als Futtterquelle.

Was? Eara, ich konnt’ doch nicht auf ihre Kleidung achten, ich war zu sehr um mein Leben besorgt! Bei meiner Seele, ich werd’ im Moment, wo Gevatter Tod seine Flügel über mich ausbreitet, seinen Schwanz um mich windet, mich zu seinem Gesicht hochhebt und mich mit seinem Feuerstrahl versengt, in dem Moment werd’ ich mich nicht gross darum kümmern, ob’s n grossen Unterschied zwischen den weiblichen und männlichen Skralen gab, Schande über mich! Ja, ’n paar hatten einfach Kindlein auf den Knien sitzen, und sahen weniger kriegerisch aus, da dacht ich nur... egal, weiter jetzt!

Ich sitz’ da also vor mich hin, umgeben von gruseligen Skralen, und mal mir aus, was die mit mir anstellen werden, sobald die Sonne aufgeht und ihre Vollmondnacht des Friedens vorbei ist, da hör’ ich ein Rummeln vom Skralberg, dann kam das näher, und dann flammte das Feuer plötzlich doppelt und dreifach auf, und kein Scherz, ich konnt ’n Gesicht darin erkennen, das grinste bösartig in die Runde.

Die Skrale jubilierten, und einer von ihnen stand auf und rief was, das wie „Darakinai“ klang. Er griff nach einem durchlöcherten, prall gefüllten Sack, welcher neben ihm stand, und zog einen riesigen Säbelfisch hervor, präsentierte ihn stolz den Anwesenden, und warf ihn dann ins Feuer, wo dieser riesige neu erschienene Feuergeist stand. Ich beobachtete, wie der silbrig glänzende Fisch vom Feuer verzehrt wurd', und der Feuergeist noch stärker grinste. Dann – mir fielen die Augen fast aus’m Kopf – krabbelte aus dem Feuer ein klitzekleiner Skral hervor, der in die Runde starrte und seine Zähne bleckte. Sofort sprang der Skral, der den Fisch ins Feuer geworfen hatte, auf, und schnappte sich das kleine Wesen.

Als wär’n Bann gebroch’n worden, standen noch viel mehr Skrale auf, und sie alle griffen in Säcke, die sie mitgebracht hatten, und warfen Fische, Eidechsen, Vyperas, behüte, sogar einen Hund glaube ich geseh’n zu haben! Also eben, sie warfen diese toten Tiere ins Feuer und kaum hatten die das Feuer berührt, krabbelten weitere kleine Skrale aus den Flammen des Feuergeists hervor und stürzten sich in die Menge. Die erwachsenen Skrale schnappten sich die Kinder, wo sie konnten, und brüllten und tanzten ausgelassen – ein Wunder, dass im Dorf auf der anderen Seite des Hügels nicht alle Bewohner senkrecht auf ihren Matratzen standen, mich hätte das Gejohle bestimmt aus dem Federn geholt. Aber dort drüben tat sich nichts, und die Skrale wurden immer wilder, und einige hauten einander und auch mir ihre Pranken auf den Rücken, und die Kindlein wurden in die Luft geworfen, und ich, mir wurde speiübel, weiter zuzusehen, darum schlich ich mich so leise ich konnte rückwärts aus der Masse raus. Die waren inzwischen so ekstatisch, die achteten nicht mal auf mich! Viele hatten sogar ihre Augen geschlossen und schmusten mit den kleinen Skralbalgen, die aus’m Feuer traten.

Dann war ich raus aus der Menge, und ich nahm meine betrunkenen Beine in die Hand, und torkelte, nein, rannte so schnell ich nur irgendwie konnte den bemoosten Hügel hinunter, zur sicheren Narne, weg von den Skralen, weg von diesem höllischen Skralberg.

Ich war kaum ’n Dutzend Schritt vom Lager entfernt, da tönte ein langgezogenes Kreischen hinter mir. Als ich mich umgedreht hatt’, sah ich nur, wie die Skralhex’, diese verfluchte Skralhexe, mir nachstarrte und auf mich zeigte, die langen Arme ausgestreckt, den Kiefer zu diesem schauderhaften Schrei geöffnet, der durch Mark und Bein ging. Und der Feuergeist, der hörte auf, Skralkindlein auf die Lande loszulassen, und hob ab von seinem Lagerfeuer, setzte sich in Bewegung und huschte über das trockene Gras hinweg auf mich zu, riesig, feurig heiss, und Zeter und Mordio spuckend. Die Skrale wollten mich doch noch als Opfergabe nutzen, ich hatt’s ja gewusst!

Doch der Feuergeist war nicht schnell genug. Es gelang dem feurigen Biest zwar, mich einzuholen – guckt ma’, die Brandnarben auf meinen Arm, die wird’ ich meiner Lebtag nimmer loswerden – aber kaum hatte der Geist mich erreicht, so war ich nah genug am rettenden Ufer, und ich warf mich in die Fluten der Narne, betend, dass ich keine Klippe treffen möge.

Ich klatschte hart auf’m kalten Wasser auf, und ging schon in den ersten paar Schwimmzügen beinahe unter – jetzt wünscht’ ich mir, einer der edlen Herren der Rietburg zu sein, die sich’s leisten können, in der Freizeit schwimmen zu gehen. Ich, ich konnt’ nur angestrengt strampeln und ich hatt’ schon zum x-ten Mal in dieser Nacht mit meim armselgen Leben abg’schlossen, da sank ich zum Grunde der Narne wie’n Stein.

Doch die Flussgeister und vielleicht sogar der grosse Arkteron höchstpersönlich müssen mir huldig gewesen sein, denn ich wurd’ wieder hochgehoben, und der Fluss spuckte mich auf der anderen Uferseite aus, weg von dem Skralreigen, weg von dem tobenden Feuergeist, weg von dieser grusligsten aller Begegnungen, die ich je in meim ganzen Leben erlebt hatte.

Nein wirklich, ich sag’s euch, die G’schichte stimmt. Was für’n Grund hätt’ ich, mir so’n verstörenden Quatsch auszudenken?! Ich hab den Skralreigen gesehen! Den Skralreigen! Und so warn’ ich euch jetzt, sodass ihr euch meine Worte für immer merken sollet: Wenn der Vollmond scheint über dem Lande Andor, so meidet den Skralberg. Meidet den Skralberg, wenn euch euer Leben lieb ist!






\newpage
\section{Die Wassermagierin, der Feuerkrieger und der Seher}



Es regnete in Strömen. Schon seit Stunden peitschten Wassermassen gegen die Fenster der Taverne zum Trunkenen Troll, doch die Tavernengemeinschaft liess sich dadurch nicht die Stimmung vermiesen. Gildas Metvorrat würde noch lange reichen, und falls man des schlechten Wetters wegen nicht nach Hause konnte, so wurde halt in der beliebtesten Taverne von Andor gleich doppelt und dreifach so lange gefeiert.

Die Tür zur Taverne öffnete sich langsam, als hätte jemand sachte dagegen gestossen.

Seltsamerweise fielen keinerlei Regentropfen durch die dunkle Öffnung ins Innere des gemütlichen Raums. Eine Hand erschien aus der Dunklen und griff an den hölzernen Rahmen der Tür, danach ein wohlbekanntes Gesicht, welches in den hell erschienenen Raum guckte – es war Jarid, die Wassermagierin aus dem fernen Danwar! Ihr sonst übliches fröhliches Lächeln fehlte, stattdessen zeigten ihre Gesichtszüge Unbehagen, ja sogar Sorge!

Sie schleppte sich in den Schankraum. Ihr selbst ging es ja prima, aber ihrem stetigen Begleiter Trieest eher weniger. Ihn sorgsam unter einer Achsel stützend, hievte sie den schwer berüsteten und vollkommen durchnässten Feuerkrieger vorsichtig über die Türschwelle und hinein in die gute Stube. Kurz flatterte sie mit ihren Fingern, und aus ihrem Haar, ihrem verzierten Umgang, ihren Schuhen, Trieests Rüstung, seinen Ohren – von überallher perlten kleine Wassertropfen hervor und schwebten, geleitet von Jarid, hinaus aus der Taverne in die finstere Sturmnacht. Mit einem dumpfen Knarren fiel die Tür hinter ihnen ins Schloss.

Der nun wieder trockene Trieest konnte sich kaum auf den Beinen halten. Er ächzte und schnaufte, doch war äusserlich keine Verletzung an ihm zu erkennen. Der fest in seiner Brust verankerte Lavastein leuchtete in einem satten Orangeton und strahlte spürbar Hitze ab. Der Stein war es wohl auch, der Trieest so sehr zu schaffen machte – die Hitze selbst schien ihm ja nie etwas anhaben zu können, doch litt er auf eine ganz andere Weise durch das Tragen dieses magischen Edelsteins.\bigskip



„Bitte, Jarid, ich kann nicht mehr lange“, flüsterte Trieest flehend aus brüchigen Lippen hervor.

Jarid blickte ihn mit einem traurigen Blick an und sprach leise: „Tri, du weisst, dass ich den Stein nicht entfernen darf, selbst, wenn ich mir sicher wäre, wie es funktioniert. Ich helfe dir ja schon beim Tragen deiner Bürde. Und ich tue, was ich kann, um die Pein zu lindern.“

„Fünf Jahre! Fünf Jahre trage ich ihn jetzt schon auf mir! Was verlangt Mutter Natur denn noch von mir?!“, stiess Trieest zwischen zusammengebissenen Zähnen hervor, ehe er wieder zum bettelnden Ton wechselte, „Nur schon fünf Minuten... fünf Minuten Ruhe und Stille... ich bitte dich...“

Jarid hielt kurz inne und schien zu überlegen, doch dann setzte sie ein strenges Gesicht auf und sagte bestimmt: „Nein. Wir wissen nicht mal, was mit dir geschehen würde, wenn ich die Verbindung kappte. Ich will den Ältesten nicht berichten gehen müssen, dass dein Prozess des Wandels nochmals von Grund auf zu beginnen hätte. Und du willst das doch auch nicht. Der Stein bleibt, wo er ist. Kümmern wir uns lieber um deine Schmerzen, da kann ich tatsächlich von Hilfe sein.“

Trieest brummelte etwas unter seinem Atem und seine langen Ohren wackelten. Dann liess er sich mühselig auf die nächstgelegene Sitzbank fallen und stöhnte schmerzerfüllt auf. Jarid quetschte sich sorgsam neben ihn.

„Gold...“, ächzte Trieest jetzt, „Hast du noch etwas von der Goldsalbe übrig?“

Jarid fischte ein kleines braunes Döslein aus einer Tasche ihres blauen Gewands und schraubte es auf. Ihr enttäuschter Blick sprach Bände: „Wasser wird es tun müssen. Viel Wasser.“

Mit diesen Worten blickte die Wassermagierin um sich, erwartend, von einem Haufen neugieriger Bauern umringt zu sein, von denen bestimmt einer ein Fass voller Wasser zu holen vermochte. Mit Erstaunen nahmen die beiden Danware jedoch wahr, dass sich noch kein Trubel gebildet hatte, ja, dass sich gar niemand gross um sie scherte. Das war aussergewöhnlich, da die beiden doch noch jedes Mal im Nu die Hauptattraktion in der Taverne zum Trunkenen Troll geworden waren, wenn sie bislang dort vorbeigeschaut hatten. Tatsächlich waren heute die meisten Besucher der Taverne aber in einer ganz anderen Ecke des Raumes um einen grossen Tisch versammelt und diskutierten angeregt.

Schnell fischte Jarid einen Trinkschlauch aus einer Falte ihres Kleids – nur noch zur Hälfte gefüllt mit abgestandenem Brunnenwasser, aber bei weitem besser als nichts – zog den Zapfen heraus und konzentrierte sich auf die klare Flüssigkeit. Mit einigen geübten Gesten, begleitet von einem dumpfen Summen und leuchtend blauem Schein, geleitete sie das Wasser aus dem Schlauch auf Trieest Brust, wo es eine Schicht über dem Lavastein bildete. Trieest seufzte auf und entspannte sich ein wenig. Doch bei der Hitze, welche der Lavastein abstrahlte, war das Wasser bald verdampft. Jarid versuchte erfolglos, die Tropfen wieder aus der Luft zu greifen. Die Tür zum Unwetter draussen wollte sie nicht mehr öffnen, wenn es sich vermeiden liess – es war sehr anstrengend, den Orkanböen Sturmwasser zu entreissen, und sie war bereits stark ausgelaugt. Gilda hatte bestimmt einige Fässer frischen, sauberen und vor allen Dingen nicht widerspenstigen Trinkwassers in ihrem Keller, welche Jarid vielleicht verwenden dürfte.

„Ich werde Gilda nach Wasser fragen, warte hier“, sprach sie zu Trieest. Dieser lachte heiser auf und bereute dies gleich wieder, als das gewohnte Ziehen in seiner Brust zurückkehrte. Wo wollte er schon hingehen in seiner jetzigen Kondition? Momentan würde es ihn nicht mal stören, von neugierigen Andori umringt zu sein. Ihre mitleidigen Blicke würden ihn bestimmt schon etwas besser fühlen lassen.\bigskip



Jarid bewegte sich auf den grossen Tisch in der Ecke der Taverne zu, um den sich die meisten Tavernengäste versammelt hatten. Sie erkannte Eara, eine Zauberin aus dem noch ferneren Hadria und gute Freundin von ihr, wie sie mit geschlossenen Augen, den Körper leicht vor- und zurückwippend, ihre Fingerspitzen über die Tischoberfläche gleiten liess, während die umliegenden Tavernengäste ihr laut Dinge zuriefen. Einige schienen auch Wetten darauf abzuschliessen, ob Earas Vorhaben gelingen würde – wenn der Handelszwerg Garz seit Taroks Niedergang auf den Kontinent zurückgekehrt wäre, so hätte man ihn jetzt bestimmt in dieser Menschenmenge gefunden, glücklich grinsend einigen Bauern ihre letzten Silberstücke aus der Tasche ziehend. Doch der Handelszwerg war nirgends zu sehen, und so blieb es vor allem den wohlhabenden Händlern vom freien Markt vorbehalten, die Wetten abzunehmen.\bigskip




Trieest beobachtete das aufgeregte Treiben von seiner Sitzbank aus angestrengt, die orange glühenden Augen fest zusammengekniffen, um möglichst scharf zu sehen. Viel mehr konnte er so jedoch auch nicht erkennen. Es schien sich um ein Spiel zu handeln, welches offensichtlich die Gemüter hochkochen liess. Doch da! Eine gebückte Gestalt löste sich aus der Menschenmenge und kam langsam auf Trieest zu. Die Person trug einen schweren dunklen Mantel und stützte sich im Gehen auf einen langen Stock – Jarid war das definitiv nicht, und auch sonst erinnerte sie Trieest an niemanden, den er kannte.

Der Feuerkrieger sah weit entfernte Dinge nur verschwommen, und seine Wahrnehmung von Farben war nicht dieselbe wie die der meisten Menschen, denen er bislang begegnet war, doch selbst er konnte erkennen, dass die Hautfarbe des nun vor ihm angekommenen Mannes sich von den hierzulande üblichen unterschied. Als dieser sich mühsam Trieest gegenüber niedergelassen hatte, fiel letzterem auf, dass der Neuankömmling gar nicht so alt war (oder zumindest nicht so alt aussah), wie Trieest zunächst gedacht hatte. Die langsamen, vorsichtigen Bewegungen und der Stock kamen daher, dass die Augen des Mannes von einer Binde bedeckt waren – er war blind!\bigskip



„Blau sind die Weiten des Meers“, sprach der fremdartiges Blinde zur Begrüssung, sehr zur Überraschung Trieests, welcher schmerzerfüllt hustete und mit krächzender Stimme antwortete: „Und Blau sind die Weiten des Himmels.“

Stille.

Trieest ergriff das Wort: „Seid... seid ihr aus Danwar?“

Der Blinde lachte, und sagte dann: „Selbst bin ich noch nie dort gewesen. Gehört habe ich allerdings viel. Von den örtlichen Gepflogenheiten über eure grossen Errungenschaften wie das Fangen der Blitze – bis hin zur Tatsache, dass kein Feuerkrieger seine Heimat verlassen sollte.“

Ein Grinsen huschte über das Gesicht des Mannes, als er seinen Satz beendete. Trieest atmete schwer. Da war viel auf einiges zu verarbeiten. Noch keiner, den er hier auf dem Festland getroffen hatte, hatte die Floskeln gekannt.

Nun gut, ganz richtig war das nicht. Reka hatte die Floskeln verwendet. Die alte Hexe hatte sich in ihrem ganzen Leben noch nie näher als zwei Meter an ein Gewässer gewagt und hasste schon das Überqueren von Brücken – dass sie je die Nebelinseln besucht hatte, hielt Trieest für sehr unwahrscheinlich. Sicherlich hatte sie damals mit Hexerei in seinen Geist geblickt und ihn auf die Art und Weise begrüsst, wie es für ihn am normalsten war.

Jedenfalls hatten die Bewahrer vom Baum der Lieder, deren Pergamentsammlung den Wissensstand der Andori ganz gut repräsentierte, noch nie von den Floskeln gehört gehabt. Ganze drei Tage lang hatten sie ihn in Anspruch genommen, um jedes Detail über Danwars Kultur aus ihm herauszuquetschen – und danach hatte sich Gända Jarid geschnappt und mit ihr dasselbe durchgezogen. Sechs volle Tage hatten sie insgesamt dadurch verloren! Aber das lag hinter ihnen. Jetzt reiss dich zusammen, Trieest, keine Abschweifungen mehr!\bigskip



Konnte es sein, dass der Blinde log und schon in Danwar gewesen war? Warum hätte er dann lügen sollen? Die Floskeln hatte er bewusst verwendet. Woher kannte er sie? Was wollte er von ihm? Und vor allem: Wie hatte ein Blinder überhaupt erst erkannt, dass da ein Danware vor ihm sass? Ein ungutes Gefühl baute sich in Trieest Magengegend auf – eine Mischung aus Verwirrung und Furcht.

Schnell blähte Trieest seine Nasenflügel auf und versuchte möglichst unauffällig, den Duft der Umgebung in sich aufzunehmen. In der schon etwas abgestandenen Tavernenluft lag nebst dem allgegenwärtigen Metgeruch und natürlich den individuellen Gerüchen der lauten Tavernengäste hinten am Tisch deutlich etwas Metallisches. Jawohl, da steckte einen Dolch im Gewand seines schweigenden Gegenübers. Andererseits war sich Trieest sicher, auch in seinem geschwächten Zustand dem Blinden körperlich weit überlegen zu sein. Gefährlich könnte es für ihn also nur werden, wenn der Mann ein Zauberer war. Und sein Stock mit dem seltsam geformten Ende sprach wahrlich dafür. Trieest verspannte sich. Dann fiel ihm auf, dass vorhin ein Hauch Eara an ihm vorbeigeweht war. Eara war die weiseste Zauberin, die der Feuerkrieger kannte (wenn auch die einzige), und wenn sie hier war und sich nicht um den Blinden sorgte, so musste er das auch nicht. Trieest entspannte sich wieder etwas.\bigskip



Den Blick auf sein Gegenüber richtend, fiel Trieest auf, dass schon eine Weile lang Schweigen zwischen ihnen beiden herrschte. Der Blinde machte keine Anstalten, das Schweigen zu brechen. Trieest erinnerte sich, dass er eine angedeutete Frage gestellt hatte. Er wollte wohl wissen, warum ein Feuerkrieger sich abseits von Danwar befand. Und Trieest hatte keine Lust, ihm das mitzuteilen. Eigentlich wollte Trieest nichts weiter, als dass Jarid mit einem Zuber voll kalten Wassers zurückkehrte und das Pochen in seiner Brust linderte. Er war müde. Aber er musste sich jetzt wohl um den geheimnisvollen Mann kümmern.

Jarid hätte in einer solchen Situation sicher mit einer Gegenfrage geantwortet – „Wie habt ihr von Danwar erfahren?“ – und bei jeder passenden Gelegenheit nachgebohrt, um an mehr Informationen zu kommen und zugleich das eigene Herausrücken von Informationen zu verhindern. Jarid war geschickt in dieser Kunst. Trieest auch, aber momentan wollte er sich nicht auf solch ein geistiges Gefecht einlassen. Stattdessen beugte er sich etwas nach vorne (er sog scharf Luft ein, als das Pochen in seiner Brust zunahm) und sprach leise, aber bestimmt: „Verzeiht bitte. Es war ein langer Tag, mir tut alles weh und ich will nichts mehr als einfach einzuschlafen und die Welt zu vergessen. Teilen sie mir doch einfach ihr Anliegen mit und dann sehen wir weiter.“

Trieest lehnte sich vorsichtig wieder zurück und hoffte, dass der Blinde seinem Appell folgen würde. Vielleicht war er ein Tulgori? Die Andori hatten erst kürzlich erfahren, dass hinter den Fahlen Bergen ein weiteres Land lag, und seitdem die Eisdämonin, die diese Gipfel gehütet hatte, erschlagen worden war, florierten die Handelsrouten. Mit den Tulgori waren auch einige Temm, kleine gescheite Wichtel, nach Andor gekommen – vielleicht gab es noch mehr humanoide Spezies von dort? Oder war sein Gegenüber doch ein Mensch? Der lange Mantel machte es schwer, viel mehr von ihm zu erkennen als seine Mundpartie.\bigskip



„Du hast immer noch die Augen des Ungeheuers, das du einst warst“, sprach der Blinde jetzt plötzlich. Das traf Trieest unerwartet. Etwas defensiv warf er zurück: „Woher wisst ihr das?“

„Ich muss es nicht mit meinen eigenen Augen erblicken, um es zu wissen. Ich bin ein Seher.“

Trieest versteifte sich wieder. Ein Seher, das hatte ihm noch gefehlt. Konversationen mit Sehern waren... kompliziert, gelinde gesagt. Oft hatten sie schon im Voraus eine Vorstellung davon, was man sagen wollte, weshalb Gespräche mit ihnen sehr rasch voranschritten und von Thema zu Thema sprangen – anstrengend für alle Beteiligten ausser dem Seher. Oder zumindest war es so beim grantigen Kord gewesen, dem Seher von Danwar.

Die Worte des Sehers hatten ihre Wirkung allerdings nicht verfehlt. Jetzt war Trieests Interesse geweckt. Jetzt wollte auch Trieest Informationen von seinem Gegenüber, nicht nur umgekehrt. Jetzt konnte verhandelt werden.\bigskip



„Wie lange...?“, begann Trieest. Eine reine Farce, falls der Blinde wirklich ein Seher war, wusste er schon im Voraus, was Trieest von ihm erfahren wollen könnte. Die Frage war vielmehr, was er dafür verlangen könnte. Gold oder Wertsachen? Da gab es hier nicht viel zu holen, Trieest besass nichts von Wert ausser seiner Rüstung und seiner Dornenklinge, die er ganz bestimmt nicht herausrücken würde. Wollte der Fremde Informationen? Der Seher schien schon so viel zu wissen, was könnte Trieest ihm da noch Neues erzählen?

„Darf ich?“, sprach der Blinde und streckte seine Hand nach dem Lavastein aus. Trieest fühlte sich erneut überrumpelt, nickte aber aufgeregt. „Vorsichtig, der Stein ist sehr h...“

Überrascht brach Trieest seine Warnung ab, als der Seher seine Hand auf den Lavastein legte. Der Edelstein leuchtete in vollster Pracht und strahlte solche Wärme aus, dass selbst Jarid etwas Abstand hätte nehmen müssen. Doch diesem blinden Seher schien die Hitze nichts auszumachen, er schien sie sogar zu geniessen. Wie das Licht so durch seine Hand schein, waren seine Fingerknochen deutlich als Schatten unter der dunklen Haut erkennbar. Doch waren da keinerlei Anzeichen von Verbrennungen oder Schmerz.

Stattdessen leuchteten die Augen des Sehers auf, in glühendem Weiss, so hell, dass es selbst durch die dunkle Augenbinde, die er trug, ein klein wenig hindurchschimmerte. Trieest versuchte, so still wie möglich zu sitzen. Dies könnte der Moment sein, der zum Abschluss seines Prozesses des Wandels führte. Heute, hier und jetzt, könnte er Informationen darüber erhalten, wie er die Bürde ablegen konnte. Jetzt musste er nur noch herausfinden, was der Seher verlangen würde, damit er seine Vision preisgab. Er lachte innerlich auf. Noch vor einigen Minuten hatte er nichts lieber gewollt, als dass dieser Mann ihn in Ruhe lassen würde, und jetzt konnte er kaum erwarten, was der Seher berichten würde.

Wie schon gesagt – bei Interaktionen mit Sehern geht alles sehr rasch voran.\bigskip



„Mutter Natur meint es gut mit dir“, schmunzelte der Seher, „Noch drei Monde, und dann wird Jarid den Stein entfernen können. Das Zeichen, welches den richtige Moment kennzeichnet, habt ihr ja bereits erraten, aber ich bestätige es gerne: Es sind deine Augen, Trieest, die schon so bald eine wunderschöne bernsteinfarbene Iris tragen werden. Damit es soweit kommt, musst du nur noch eine weitere – eine letzte – Prüfung deiner Entscheidungskraft durchstehen.“

Eine Vielzahl von Gefühlen schoss durch Trieest – Erleichterung, Freude und Anspannung. Vergessen war seine Müdigkeit, vergessen waren für einen kurzen Moment sogar seine Schmerzen. Drei Monde! Das war nichts im Vergleich dazu, was hinter ihm lag. Und er Tor hatte sich gefürchtet, dass es doppelt und dreifach so lange dauern hätte dauern können! Eine letzte Prüfung lag noch vor ihm, hatte der Seher behauptet. Worin die wohl bestand?

„Ich kann dir nicht sagen, worin deine Prüfung besteht, aber ich kann dir sagen, wo sie stattfinden wird. Ich sah es vor meinem inneren Auge, so deutlich, wie ich nur selten eine Vision hatte – die Insel Narkon ist es, der du dich stellen musst. Ziehe von Süden her über die Mauerberge. Dort wird der letzte Test auf dich warten. Dann wirst du wissen, worin diese Aufgabe besteht.“

Jetzt war Trieest euphorisch. Er hatte ein Ziel! Narkon! Eine geheimnisvolle Nebelinsel, über die nicht viel bekannt war. Das Ende seiner Tortur war nahe. Danach... er hielt inne. Ja, es gab endlich Aussicht auf ein „danach“ für ihn! Viel von der Welt hatten Jarid und er inzwischen gesehen, vielleicht war es ja nach dem Ablegen seiner Bürde Zeit, sich zur Ruhe zu setzen. Wollten sie nach Danwar zurück? Ihre Freunde und Familie wiedersehen? Oder sich lieber fernhalten, die alten Wunden nicht wieder aufreissen, und ein neues Leben in einem fernen Land beginnen? Er lachte. Egal, wozu sie sich entscheiden würde, es würde grossartig werden. Bald, bald würde er frei sein!

Als Trieest seine Aufmerksamkeit wieder zum Seher wandte, fragte er verwirrt: „Wartet... warum teilt ihr mir das schon alles mit? Was wollt ihr denn für eine Belohnung?“

Aufrichtig lächelnd, sprach der blinde Seher: „Nicht alle Seher verlangen eine Bezahlung. Der Seher von Danwar – ist es immer noch der grantige Kord? – sollte es wahrlich besser wissen!“\bigskip



Trieest lachte auf. Es war ein wundervoller Tag. Sobald dieser Sturm sich gelegt hatte, konnten Jarid und er in den Wachsamen Wald reisen, um von dort aus ein Händlerschiff nach Sturmtal zu nehmen. Die Taren waren freundlich und ihre Insel lag ganz nahe an Werftheim, was nun wirklich ein Streifenmardersprung von der Silbermine entfernt lag, und vor dort aus würden sie die Mauerberge überqueren, und dann... dann würden sie weitersehen. Trieest war zuversichtlich. Er hatte wieder ein Ziel, einen Sinn. Er war wunschlos glücklich.

Plötzlich flackerte ein Ansatz von Unsicherheit über die Mundpartie des Blinden. Dann hatte er sich wieder unter Kontrolle, drehte sich mit einem breiten Lächeln auf dem Gesicht um und sprach: „Grün sind die Wogen der We...“

„Spar dir das Gefasel“, unterbrach Jarid ihn kalt. Die Wassermagierin war endlich zurückgekehrt, ein grosses Fass in den Armen. Wie lange hatte sie bereits hinter dem Seher gestanden? Als Trieest ihre vor Wut blitzenden Augen sah, wusste er, dass irgendetwas nicht in Ordnung war. Seine Hoffnung schwand.

Jarid ignorierte den Seher und wandte sich direkt an Trieest: „Glaube ihm kein Wort! Das ist ein Lügner der schlimmsten Sorte!“

Trieests Euphorie brach in sich zusammen wie ein Gor, dem man zünftig eins auf den Schädel gehaut hatte. Er kannte Jarid gut genug, um zu wissen, dass sie solche Anschuldigungen nicht leichtfertig machen würde. Und als sein Blick zurück auf den Seher fiel, sah er es in dessen Mimik, unter der Kapuze – ein kurz aufblitzender Anflug von Ärger und Verdruss, aber auch Enttäuschung. Ein Anblick, der Jarids Aussagen Recht zusprach. Der Seher war ein Lügner!

Der Blinde gab sich vergeblich Mühe, seine gute Miene wieder aufzusetzen. Sein Kopf ratterte bestimmt auf Hochtouren, versuchend, eine Möglichkeit zu finden, die Situation wieder zu seinen Gunsten zu wenden. Erfolglos, offenbar, denn der Seher erhob sich zackig, verbeugte sich – fast schon spöttisch – vor der Danware und stürmte überraschend schnell für einen Blinden und ohne zu zögern aus der Taverne hinaus in stürmische, nasse, kalte Nacht. Wenn er es nicht besser gewusst hätte, hätte Trieest geglaubt, eine Träne über das Gesicht des Sehers geflossen gesehen zu haben, ehe die Tür hinter ihm zuschlug.\bigskip



Jarid setzte sich vorsichtig neben Trieest und legte ihm eine Hand auf die Schulter. Sie sagte nichts. Trieest blieb ebenfalls ruhig. Still teilten sie gemeinsam seine Trauer, den Verlust einer Hoffnung, die sich so schnell wieder verzogen hatte, wie sie gekommen war. Dann nahm Jarid all ihre Kräfte zusammen und öffnete das Wasserfass, welches ihr Gilda netterweise überlassen hatte. Vorsichtig, um ihre magische Kapazität nicht zu überstrapazieren, führte sie einen dünnen, blau schimmernden Wassersstrang aus dem Fass zu Trieests Lavastein und flocht sorgsam ein abschirmendes Netz um den Edelstein in Trieests Brust. Dieser seufzte auf.

„Woher wusstest du... “ setzte der Feuerkrieger an, doch ein Hustenanfall schüttelte ihn – beinahe wäre Jarids Wassergeflecht gebrochen! – und er schloss den Mund. Jarid hatte ihn verstanden.

In einer etwas zu hohen Stimmlage antwortete Jarid vorsichtig: „Tri, ich habe dir nicht die ganze Wahrheit gesagt.“

Trieest blickte sie stumm an, die orange glühenden Augen ausdruckslos, doch fühlte Jarid, wie er sich wieder verspannte. Seine Hände ballten sich zu Fäusten.

Ihre nächsten Worte sehr sorgfältig wählend, antwortete Jarid:

„Bevor wir Danwar... verliessen, brachten mich die Ältesten zur roten Grotte.“

Trieest blickte Jarid überrascht an. Diese fuhr stockend fort:

„Dort wurde mir... uns... eine Prophezeiung auf den Weg gegeben. Ich... ich habe dir nichts davon gesagt, weil... weil... ich durfte nicht. Ich darf nicht. Sie sagten, dass du... dass, wenn ich es dir sage... Tri...“

Trieest blickte Jarid noch überraschter an. Es kam nicht oft vor, dass sie die Fassung verlor. Doch jetzt waren Jarids Augen feucht, und auch wenn Trieest wusste, dass sie den Tränen nicht erlauben würde, zu fliessen, so war er doch besorgt. Beim Barte des Warx, was hatte das Orakel Jarid denn gesagt, dass sie solch eine Reaktion zeigte?

Jarid nahm einen tiefen Atemzug, dann hatte sie sich wieder etwas gefasst: „Ich weiss, wo du deinen Prozess des Wandels beenden wirst – und es ist nicht Narkon.“

\begin{center}
    Weiter geht es in \hypref{Der Feuerkrieger und die Wassermagierin (2023)}.
\end{center}




\newpage
\section{Eine mysteriöse Tafel}

Rund um die Taverne zum Trunkenen Troll tobte ein Sturm, wie man in Jahren keinen gesehen hatte. Die Tavernengäste, auf der Suche nach einer Beschäftigung, hatten sich einmal mehr den mysteriösen Tafeln zugewandt, die Gilda bei einer seltsamen Begegnung erhalten hatte.

Nicht einmal der alte Leander, der einst behauptet hatte, alle Sprachen der bekannten Welt zu kennen, hatte es geschafft, die Nachrichten zu entziffern. Der mürrische Blinde hatte zwar behauptet, das läge daran, dass er nicht gut ertasten könne, was für Zeichen auf der Tafel eingeritzt waren, und wenn er das Bild nur vor Augen hätte, würde ihm die Bedeutung sofort klar – er sei schliesslich ein Seher. Die meisten hatte diese Aussagen allerdings als Humbug abgetan, als Ausreden eines Mannes, der es zu einer solchen Meisterschaft in einer Disziplin gebracht hatte, dass er schlicht verlernt hatte, zuzugeben, etwas nicht zu können.

Ebenjener Leander wandte sich soeben ab von Earas erneutem Versuch, die Tafel mithilfe von Heller Magie zu entschlüsseln – was für eine Närrin sie war, so etwas überhaupt mit Zauberei zu probieren. Mit dem trotzigen Auswurf „Ich weiss doch schon, ob es ihr gelingen wird. Bin ja ein Seher!“ entfernte er sich von der immer noch lautstark mit Eara miteifernden Menge. Leander hatte eine interessantere Präsenz in der Ecke des Raumes gespürt. Die Tavernengäste beachteten seinen Abgang nicht – offenbar würden sie ihn nicht vermissen. Innerlich grinste Leander aufgeregt. Er liebte es, den mürrischen Alten zu spielen. Und noch wichtiger: Das Schicksal hatte ihm soeben eine Möglichkeit geschenkt, die Rettung seines Bruders in Gange zu setzen.



Eara liess ihre Fingerspitzen mit geschlossenen Augen über die Oberfläche einer der eigenartigen Schriftstücke gleiten und wippte ihren Körper leicht vor- und zurück, während sie sich einen Zauberspruch der Erkenntnis und der Übersetzung in Erinnerung rief und gleich mehrmals hintereinander flüsterte.

Gilda beobachtete Earas Tun amüsiert – sie glaubte nicht, dass es beim fünften Versuch anders enden würde, aber Eara war ehrgeizig und nur sehr schwer von etwas abzubringen, wenn sie es sich einmal in den Kopf gesetzt hatte. Zudem lenkten ihre Versuche, das Schriftstück auf magische Art und Weise zu übersetzen, die Gäste vom tobenden Sturm ab, und was wäre Gilda für eine Tavernenbesitzerin, wenn sie sich nicht um die Freude ihrer Gäste gekümmert hätte?

Plötzlich wurde Gilda von hinten auf die Schulter getippt. Sie blickte sich um und war überrascht, Jarid dort zu erblicken – sie hatte die Danwarerin gar nicht hereinkommen sehen. Nun gut, wenn irgendjemand in diesem Unwetter umherreisen konnte, dann war es wohl eine Wassermagierin. Ihrer Anwesenheit nach befand sich Trieest wahrscheinlich auch schon in der guten Stube, und das hiess...

Ehe Jarid ihre Bitte überhaupt stellen konnte, hatte Gilda Zwei und Zwei zusammengezählt und fragte: „Ich hab‘ mehrere Fässer sauberen Wassers im Keller, wie viele braucht Trieest?“

Dankbarkeit blickte in Jarids Gesicht auf, als sie antwortete: „Mehr als eines ist wirklich nicht nötig. Und bemüh‘ dich nicht, ich hole es mir selbst gleich.“

Sie zögerte kurz und beugte sich noch näher zu Gilda, um leise zu wispern: „So sage mal, was macht Eara da?“

Gilda grinste und sprach: „Sie verzweifelt am Übersetzen einer unübersetzbaren Schiefertafel. Sobald es Trieest etwas besser geht, könntet ihr euch auch mal daran versuchen. Wer weiss, vielleicht ist das ja gerade die glorreiche Aufgabe, die seine ‚Wandlung‘ vollenden kann. Na?“

Jetzt war es an Jarid, kurz aufzulachen: „Das sind doch keine unübersetzbaren Tafeln, das ist kinderleicht! Solche benutzten wir schon während unserer Schulzeit in Danwar – du verstehst, in einer so feurigen Welt arbeitet man lieber nicht mit Schriftrollen aus Papier. Die Zeichen sind Danwarische Stichrunen, die kann ich besser lesen als eure andorischen Lettern! Schau mal, da im oberen Teil steht zum Beispiel ‚S.134: Jeder wackere Bauer vertilgt bequem zwei Pfund Gorfleisch.‘ Zugegeben, dieser Satz wirft mehr Fragen auf, als er beantwortet. Ich komme später gerne nochmals darauf zurück, für den Moment muss ich mich allerdings leider entschuldigen – Trieest leidet.“

Mit diesen Worten verzog sich Jarid in den Keller der Taverne und liess eine überraschte Gilda zurück. Eara hatte ihr Tun unterbrochen und starrte der Wassermagierin mit grossen Augen nach. Offenbar hatte sie nicht mit einer so einfachen Lösung für ein so lange ungeklärtes Rätsel gerechnet. Hoffentlich nahm sie es nicht persönlich, dass ihre Zauberei daran gescheitert war.

Als Eara sich gleich wieder der Tafel zuwandte und mit einem freudigen Lächeln auf dem Gesicht ein kleines Notizbüchlein und einen Stift hervorzog, wurde Gilda klar, dass die Zauberin das Ganze eindeutig nicht persönlich nahm. Jetzt wollte sie aber erst recht den auf den Tafeln stehenden Text entschlüsseln. Und wenn sie ehrlich mit sich war, war Gilda auch schon ziemlich gespannt darauf, was auf den Tafeln stand.








\newpage
\section{Endlich entschlüsselt}



„... und dieses Zeichen hier sollte ein ‚B’ darstellen.“

„Warte, das sieht eher wie ein ‚D’ aus.“

„Du hast recht! Also, dann lautet dieses Wort F – E – U – E – R – D – ?“

„Zwei Silberstücke auf ‚Ä’!“

„Drei Silberstücke für ‚R’!“

„Der nächste Buchstabe ist ein ‚R’!“

„Feuerdrache! Das Wort ist ‚Feuerdrache’!“\bigskip



Die andorischen Bauern und Händler in der Taverne zum Trunkenen Troll waren Feuer und Flamme für das Entschlüsseln der rätselhaften Tafel mit der danwarischen Schrift. Sie konnten auch nicht viel anderes tun, solange es draussen weiterhin zu heftig stürmte, um sicher nach Hause gelangen zu können. Jetzt, wo Jarid, die weise Wassermagierin aus Danwar, ihnen den entscheidenden Hinweis gegeben hatte, war das Dekodieren des Texts auch gar nicht kompliziert. Unterstützt durch mehr oder weniger hilfreiche Kommentare von Chada und Thorn, notierte Eara Buchstaben um Buchstaben in einem kleinen blau gebundenen Büchlein, welches sie stets bei sich trug. Gilda guckte ihr immer wieder mal über die Schulter und rief einen weiteren Buchstaben aus, die Spannung zwischen den Lettern natürlich so sehr in die Länge ziehend, dass die Tavernenbesucher auch genügend Zeit hatten, einen kleinen Teil ihres hart erarbeiteten Geldes darauf zu verwetten, welcher Buchstabe als nächstes kommen würde. Dabei würden sich die meisten von ihnen nicht gross für den Text interessieren, wenn er in andorischen Buchstaben geschrieben worden wäre. Aber die Stimmung war grossartig, und das war es, was Gilda zu erreichen erhoffte.\bigskip



Triumphierend verkündigte Gilda jetzt: „Die Überschrift lautet \textbf{Die Agren, der Schildzwerg und der Feuerdrache!“}

„Was ist das denn, etwa ein Märchen?“

„Die Form des Titels deutet auf ein Märchen hin. Aber die alleine hat noch nichts über den Text auszusagen.“

„Wer würde ein Märchen in danwarischen Runen in eine Steintafel ritzen und diese danach durch Gildas Fenster werfen?“

„Und vergiss nicht dieses verwaschene Bild nebendran. Soll das einen Hornbären darstellen?“

„Diese Angelegenheit schreit doch förmlich nach einem Spässchen der Waldgeister!“

„Ich muss zugeben, ich hatte mir etwas Besseres als eine Gutenachtgeschichte erhofft.“

„Ach Gord, so schweig stille! Du musst ja nicht zuhören, wenn du nicht magst.“

„Meine Mutter hat mir die Märchen jeweils nicht vorgelesen, sondern vorgesungen. Gilda, willst du das versuchen?“

„Ha! Da ist der gute Grenolin aber gar nicht glücklich, dass du nicht zuerst an ihn gedacht hast.“

„Was habt ihr denn? Laute spielen kann unser Greno doch prima!“

„Eara, brauchst du noch lange? Gilda, wie weit ist Eara schon?“

„Die Arbeit ist schon fast vollendet“, sprach Eara, ohne von ihrem Büchlein aufzublicken.

„Dieses Zeichen hier sollte doch ein ‚T’ sein“, warf Thorn ein. Eara funkelte ihn an, korrigierte den Buchstaben, kritzelte einen letzten entschlüsselten Satz in ihr Büchlein und übergab dieses dann breit lächelnd an Thorn: „Warum tust du uns nicht die Ehre, die Geschichte vorzutragen?“

Thorn zögerte zunächst, ergriff das Büchlein dann aber und verkündete:\bigskip



\textit{So vernehmet diesen Tatsachenbericht aus der Zeit des Unterirdischen Krieges, einer unsicheren, gesetzlosen Zeit. Dies ist die Legende von der Agren, dem Schildzwerg und dem Feuerdrachen.}\bigskip



„Zu schade, dass Kram in der Mine zu tun hat. Der mochte Überlieferungen aus der Vorzeit wie sonst keiner“, warf ein Schildzwerg ein.

„Lauter“, rief ein Bauer aus einer hinteren Reihe, und Thorn wiederholte mit etwas festerer Stimme:\bigskip



\textbf{\textit{Die Agren, der Schildzwerg und der Feuerdrache}}\bigskip



\textit{„DRACHEE!“}

\textit{Der verzweifelte Schrei hallte durch die Täler des Grauen Gebirges und schreckte allerlei Kleingetier auf, welches schleunigst das Weite suchte, selbst wenn es die Warnung nicht verstehen konnte. Die wenigen Agren, die sich ausserhalb ihrer sicheren Höhlen aufhielten, schreckten auf. Sie rannten allerdings nicht blindlings davon, sondern richteten ihre Augen furchtsam gen Himmel, blickten nach oben, suchten den kleinen schwarzen Punkt, welcher Feuer und Tod ankündigen würde.}

\textit{Gouri, eine vergleichsweise junge Agren (sie hatte doch erst kürzlich ihren vierzigsten Jahrestag gefeiert), versuchte ...}

[Siehe \hypref{Die Agren, der Schildzwerg und der Feuerdrache (2020)}.]

\textit{... einen Anflug kindlicher Freude. Sie würde diesen Zwerg zu den Hallen der Zwergenkönige begleiten! Sie würde auf ein Abenteuer gehen!}

\textit{Erst als die beiden die Alte Zwergenstrasse erreichten und eine weitere Höhle passierten, fiel Gouri ein, dass die anderen Agren wohl gerade umkamen vor Sorge um sie.}\bigskip







„Hier ist der Bericht zu Ende“, meinte Thorn erstaunt. „Doch die Geschichte scheint noch nicht so, als wäre sie zu Ende.“

Eara nickte: „Wenn das ein Tatsachenbericht darstellen soll, so wurde er bestimmt einige Zeit nach diesen Vorfällen niedergeschrieben, und da würde es nur sinnig sein, dass auch der Rest davon irgendwo notiert wurde. Aber wo?“

Gilda meldete sich zu Wort: „Vielleicht finde ich ja noch mehr Tafeln, wenn ich mehr Butterbrote backen würde und in den südlichen Wald brächte.“

„Sind wir sicher, dass die Tafeln von da kommen? Das sind Jarids Urteil nach eindeutig danwarische Runen.“

„Aber die Geschichte stammt aus dem Grauen Gebirge! Weiter entfernt von Danwar geht’s doch kaum!“

Die Helden sahen einander nachdenklich an.

Die anderen Tavernengäste tratschten indes munter weiter: „Ein Märchen war das ja wahrlich nicht. Diese Geschichte hat ja nicht mal irgendeine moralische Lektion dahinter!“


\newpage
\section{Reisetagebuch von Wrort, dem reisenden Temm}


\textit{Wenn ihr Lust habt, die Fortsetzung dieser Geschichte zu hören, so bin ich gerne bereit, sie zu erzählen – aber so unordentlich, wie ich bin, scheine ich die nächsten Tafeln schlicht und ergreifend verlegt zu haben! :oops:}

\textit{Das einzige, was ich noch in meiner Schriftensammlung auftreiben konnte, war ein Bericht von Wrort, dem reisenden Temm. Vielleicht kann dieser ja bei der Suche helfen?}\bigskip



Da stand ich nun also endlich in Andor, dem legendären Drachenland! Ich hatte eine Höhle erreicht, in der ein netter Fährtenleser mit seinem sprechenden Raben lebte, und teilte eine einfache Mahlzeit mit ihm (grösstenteils Apfelnüsse, diese wachsen in diesem Reich einfach überall). Er erzählte von allerlei Geschichten aus diesen Landen jenseits des Kuolema und fachte damit meinen Entdeckergeist nur noch mehr an. So brach ich am nächsten Morgen nach einer grossen Mütze tiefen Schlafs freudig in Richtung der von König Brandur und seiner Schar erbauten Burg auf. Dieses architektonische Wunderwerk wollte ich mir genauer ansehen. Ich hielt mich auf meinem Weg \textbf{so nahe wie möglich} am Kuolema, da ich mich nahe der Berge sicherer fühlte – die Kreaturen, von denen es in Andor den Gerüchten nach nur so wimmeln sollte, kamen ja bekanntlich vor allem aus dem Osten.

Schon nach zwei Stunden stand ich vor dem majestätischen Gebäude.

Ich notierte mir meine momentane Position.\bigskip



König Brandur war ein gütiger Mann, der den Worten eines Fremden rasch Glauben schenkte, und so gewährte er mir Einlass in seine Burg und liess mich sogar von einem Wachmann herumführen. Dieser fragte mich, aus welchen Landen ich den stammte, und ich berichtete ihm von Tulgor. Seine Augen leuchteten, als er von unseren grossen Kränen und Baumaschinen hörte. Dann berichtete er, dass am anderen Ufer der Narne, im Wachsamen Walde, der Orden der „Bewahrer vom Baum der Lieder“ herrschte, der meine Berichte über Tulgor liebend gerne in sein Archiv aufnehmen würde. In den Worten des Wachmanns: „Diesen eitlen Pinselnschwingern läuft doch schon doppelt und dreifach das Wasser im Munde zusammen, wenn sie nur hören, dass hinter dem Fahlen Gebirge etwas anderes als eine leere Einöde liegt.“

Ich versprach dem Wachmann, es mir zu überlegen, und brach wieder auf. Die Rietburg verliess ich durch das südliche Tor. So reiste ich fünf Stunden lang, vom ewigen Feuer vor der Burg aus am düsteren Krähenstamm vorbei bis zum Freien Markt der Händler der Andori und schliesslich zu einer kleinen, hell erleuchteten Taverne, wo ich mir mit einigen Silbermünzen eine Übernachtungsgelegenheit ersteigern konnte. Ich plauderte mit einigen Tavernengästen und auch sie rieten mir grösstenteils, zum Wachsamen Walde weiterzureisen. Doch vorher wollte ich noch etwas mehr vom Rietland sehen und das Leben der hiesigen Bauern hautnah miterleben.

Das östliche Rietland hatte ich noch gar nicht besucht, darum heuerte ich zwei Andori an, mich dorthin zu begleiten und mir auf dem Weg die Dunklen Kreaturen vom Leibe zu halten. Es waren beide einfache, aber gute und tapfere Menschen, und ich weiss nicht, wie weit ich ohne sie gekommen wäre. Kaum hatten wir die nächstgelegene Brücke überquert, so wurden wir auch gleich von einem waschechten Troll angegriffen! Ich hatte noch die in meinem Leben so ein furchteinflössendes Wesen gesehen. Die beiden Andori hielten mir das Ungetüm vom Leibe, während ich zum angrenzenden Bauerngehöft floh. Dort wurden ich und die beiden Andori gut versorgt. Wir unterstützen den Bauern als Gegenleistung bei ihrer Arbeit, so gut wir konnten (was in meinem Fall nicht sonderlich gut war, aber was soll’s). Nach dem Abendessen – leckere andorische Butterbrote – trugen mich meine kleinen Beinchen noch gut eine Stunde lang eine kleine Strecke der Narne entlang gen Süden. Dort setzte ich mich an eine Böschung und meditierte in der Stille der Nacht.

Ich notierte mir meine momentane Position.\bigskip



Beim nächsten Sonnenaufgang fanden mich die beiden Andori frierend an diesem Ort, und ich musste den beiden versprechen, nie mehr einfach kommentarlos zu verschwinden. Sie hatten sich wohl Sorgen um mich gemacht. Wie nett.

Wir drei nutzten an jenem Tage die 7 sonnigen Stunden des Tages zu ihrer Völle aus und reisten zum Wachsamen Wald und sogar bis in den Wachsamen Wald hinein, auf den sagenumwobenen Baum der Lieder zu. Auf unserem Weg überquerten wir eine auf Pfählen gebaute Brücke, die sogenannte Bogenbrücke. Davon würde der Hohe Architekt Tulgors bestimmt gerne hören.

Ich war froh, als wir die nebligen Gebieten nahe der Narne und des Likko verliessen, da sich darin jederzeit weitere Dunkle Kreaturen verstecken konnten.

Ganz bis zum legendären Baum der Lieder schafften wir es an diesem Tag nicht mehr, aber ich konnte den mächtigen Stamm durch das Fernrohr, welches einer meiner beiden Begleiter stets bei sich trug, bereits ganz gut erblicken, als wir nach 7 Stunden Wanderung auf die ersten Bewahrer vom Baum der Lieder stiessen.

Diese Bewahrer waren misstrauischer als die Bewohner der Rietburg und konnten sich nur langsam davon überzeugen lassen, ihre Bögen zu senken. Den beiden Andori vertrauten sie natürlich, aber mit mir und meinen magischen Kräften, die ich ihnen demonstrierte, war es etwas anders. Offenbar waren Zaubertricks in diesen Landen nicht mehr so gut angesehen, seitdem ein gewisser Dunkler Magier hier sein Unwesen getrieben hatte. Gut zu wissen.

Schlussendlich zeigten die Bewahrer uns aber trotzdem den kleinen Lagerplatz, den sie an dieser Stelle gebaut hatten.

Ich notierte mir meine momentane Position.\bigskip



Dann legten wir, das heisst ich, die beiden angeheuerten Andori und die beiden Bewahrer, die uns gefunden hatten, uns schlafen.

Am nächsten Tag dann trafen wir endlich beim Baum der Lieder ein. Der Oberste Priester, Melkart war sein Name, war gekleidet in ein edles weissen Gewand und tatsächlich komplett aus dem Häuschen, von einer neuen Welt hinter dem Kuolema zu hören. Er quetschte mich bei allerlei gutem Essen über die Geographie, Kultur und dergleichen meiner Heimat aus. Die beiden Andori kehrten bald darauf ins westliche Rietland zurück und die beiden Bewahrer begaben sich wieder auf ihren Wachposten

Und ich, ich werde nach einigen Tagen ergiebiger Gespräche mit Melkart und den seinen fröhlich weiterreisen. Vielleicht kehre ich später einmal zurück, um von meinen weiteren Erlebnissen zu berichten.




\newpage
\section{Ein danwarischer Epilog}

In einem gemütlichen Zimmer der Taverne hatte Trieest sich auf dem Bett breitgemacht, während Jarid neben ihm sass, immer wieder mal ihren Arm schwenkte und damit eine neue Schicht frischen Wassers über den Lavastein zog. Langsam beruhigte sich Trieests Brust. Jarids Behandlung tat ihm gut. Seine Aufgebrachtheit schwellte mit der aufsteigenden Müdigkeit ab. Das war ein anstrengendes Ende eines anstrengenden Tags gewesen. Jarid hatte vor der Verbannung einen Orakelspruch erhalten, das war die einzige neue Information, die wirklich zählte. Wie lange hätte sie es noch vor ihm verheimlichen wollen? Was war ihr nur gesagt worden, das sie selbst in Gedanken daran so sehr durcheinanderbringen konnte? Natürlich wusste Trieest, dass sie nicht darüber sprechen sollte, ja, nicht darüber sprechen durfte, dass etwas Schreckliches eintreten würde, wenn sie zuviel verriet. Aber ein kleiner Schubs in die richtige Richtung würde sicherlich nicht schaden können. Er wagte sich vor:

„Wie lange wird die Bürde noch mein sein?“

Stille.

„Ist die letzte Prüfung schwer zu meistern?“

Stille.

„Wenn du entscheiden dürftest, wohin wir als nächstes reisen, was würdest du vorschlagen?“

Stille.

„Jarid, so sprich doch!“

Stille.

„Weisst du, wir könnten dennoch nach Narkon aufbrechen.“

„Ich will nichts tun, wozu dich ein lügnerischer Zukunftswisser überreden wollte. Und still jetzt, dein Körper muss genesen.“

„Ich fühle mich schon viel besser, dir zum Dank – nein, bitte nicht aufhören! Ich meine nur: Wenn dieser Blinde wirklich ein Seher war, so werden wir letztendlich ohnehin das tun, zu was er uns bringen wollte. So könnte er etwa deine Reaktion vorhergesehen und den wütenden Abgang gespielt haben, damit wir gerade nicht nach Narkon gehen, obwohl wir das eigentlich sollten.“

„Kein Seher verfügt über derartige Allwissenheit. Das wäre zu viel für einen menschlichen Geist.“

„Nicht mal der grantige Kord?“

„Insbesondere nicht der grantige Kord.“

Stille.

„Wohin willst du als nächstes ziehen, Tri?“

„Könnten wir uns mit Merrik absprechen? Vielleicht gibt es noch ein uns unbekanntes Reich, welches Kartographierung bedürfen könnte.“

„Gerüchten zufolge ist Merrik leider kürzlich in den Norden geschifft.“

„Etwa nach Narkon??“

Seufzer.

„Ich bin ja schon still.“

Stille.

„Du, Jarid?“

„Hm?“

„Nach den Berichten der Temm, die ich vernommen habe, klingt Tulgor ziemlich spannend.“

„Dann wird also Tulgor unser nächstes Ziel. Ich werde morgen zum Baum der Lieder springen, um mich von Melkart bei der genauen Wahl des Reisepfads beraten zu lassen.“

„Wundervoll!“

Stille.

„Das kommt schon, Tri. Bleib tapfer.“

„Danke. Du musst den Orakelspruch nicht für dich behalten, das weisst du. Eara könnte dir sicherlich beim Interpretieren helfen. Oder die alte Reka. Die können so was.“

„Ich weiss. Danke.“

Stille.

„Gelb ist die Sonne in der Senke.“ Gute Nacht.

„Und Gelb sind die Blumen in der Blüte.“ Danke gleichfalls.

\begin{center}
    Weiter geht es in \hypref{Der Feuerkrieger und die Wassermagierin (2023)}.
\end{center}

